\chapter{Application Case Study: AI-Inference Serving}\label{application-case-study-ai-inference-serving}

The thesis's abstraction is not confined to classical matching markets; it instantiates a
contemporary systems problem. This chapter casts request routing in an AI-inference serving
system as online matching and shows the framework recovers the field's established results.
It is presented, deliberately, as a \textbf{case study, not a novelty claim} --- the systems facts
below are known, and the with-predictions vocabulary is a re-labeling of them (a decision
justified by the literature review of Chapter 8 and the negative probe summarized in §10.2).

\section{The serving instantiation}\label{the-serving-instantiation}

We map serving to online \(b\)-matching: the offline resources are model replicas or cache
shards, each with a \textbf{capacity} \(c\) (it can serve up to \(c\) concurrent requests, hence
\(b\)-matching rather than \(1\)-matching); arrivals are requests drawn from a non-uniform,
bursty traffic distribution over request types; an edge is a capability or cache affinity;
and \textbf{goodput} --- the fraction of requests served --- is the competitive ratio against the
\(b\)-matching optimum. We instantiate this on three real traces (Wikipedia pageviews, an
Azure LLM inference trace, and the Mooncake prefix-cache trace) and across four serving
concerns.

\begin{itemize}
\tightlist
\item
  \textbf{Capacity as robustness.} Increasing the capacity \(c\) smooths the effect of a bad
  prediction: a capacity-aware baseline stays safe, while \emph{blindly} trusting a traffic
  forecast degrades \emph{further} as capacity grows. Capacity is thus a \emph{substitute} for
  algorithmic robustness --- the systems analogue of the robustness-insurance thesis.
\item
  \textbf{Forecasts vs live load.} Under dynamic service times (requests hold a slot for a real
  duration, released event-by-event), a live-load signal beats a stale traffic forecast --- a
  reactive load balancer outperforms a forecast-following router when the forecast has aged.
\item
  \textbf{Cache-affinity routing.} For prefix-cache-aware routing, a \emph{stable} affinity router
  beats a reactive one --- the reverse of the traffic-forecast case, because cache locality
  rewards persistence.
\end{itemize}

Each of these recovers an established systems result cleanly, demonstrating that the
framework instantiates the problem faithfully.

\section{A probe for a genuinely new result, and its negative}\label{a-probe-for-a-genuinely-new-result-and-its-negative}

Because the literature suggested the with-predictions lens might yield a \emph{new} actionable
serving result on a tail objective, we probed it (the full account is in Chapter 8). On an
SLO/tail objective --- protecting a tight-SLO class of requests under bursty load --- a
non-predictive policy (static headroom, or a reactive-adaptive reservation based on observed
load) matches a clairvoyant oracle to within \(\le 3\%\) across every regime swept; a trivial
static reservation of one slot even beats the oracle in the moderate regime. Two reasons,
both robust: protecting a tight-SLO minority needs only a small static headroom, no forecast;
and bursts are persistent enough that reacting to observed load is nearly as good as
forecasting it. The tail objective is forgiving too --- a third face of the thesis's wall,
after throughput (Chapters 4--7) and predictor-learning (Chapter 8). We found no natural
regime where foresight helps, so serving remains a case study.

\section{Chapter summary}\label{chapter-summary}

The serving instantiation shows the framework reaches a live systems problem and recovers its
established results --- capacity as a robustness substitute, live load over stale forecasts,
stability for cache locality --- and that even a tail objective does not open a genuine
with-predictions win. The application is a faithful case study of the thesis, not an
independent contribution.
